\section{Task 3: Hemisphere Dwell, Crossing Point, and Apoapsis Placement}
\label{sec:task3}

\begin{newpar}
\usertxt{Because $\Omega$ only shifts the groundtrack in longitude (\cref{sec:task2}), the results from this first part will only depend on $i$, $e$, $T$, and $\omega$. To simplify, we define $\zeta = \omega + \nu$, which is the angle from the ascending node to the spacecraft. We can rotate $\vect{r}_P$ to get the ECI position:}
\begin{equation}
\begin{split}
\vect{r} &= \mathbf{R}_3(\Omega)\,\mathbf{R}_1(i)\,\mathbf{R}_3(\omega)\,\vect{r}_P \\
&= r\begin{bmatrix} \cos\Omega\cos\zeta - \sin\Omega\cos i\sin\zeta \\ \sin\Omega\cos\zeta + \cos\Omega\cos i\sin\zeta \\ \sin i\sin\zeta \end{bmatrix}.
\end{split}
\label{eq:t3-r}
\end{equation}
\usertxt{And because ECI and ECEF both share the Z component, the geocentric latitude of any point in the orbit is defined as:}
\begin{equation}
\sin\phi = \frac{z_e}{r} = \frac{z}{r} = \sin i \,\sin\zeta.
\label{eq:t3-lat}
\end{equation}

\begin{figure}[H]
\centering
% Geocentric latitude phi at a point on the orbit (schematic, not to scale).
\begin{tikzpicture}[>=Latex, scale=0.7, font=\small]
    \coordinate (O) at (0,0);
    \coordinate (F) at (2.3,-0.55); % foot of the perpendicular in the equatorial plane
    \coordinate (S) at (2.3,3.34);  % spacecraft

    % Earth and equatorial plane
    \fill[gray!20] (O) circle (0.7);
    \draw[gray] (O) circle (0.7);
    \draw[dashed, gray] (O) ellipse (3.6 and 1.1);
    \node[gray, align=center] at (-5.0,0.2) {equatorial\\plane};

    % orbit
    \draw[thick, hwblue, rotate=60] (O) ellipse (4.1 and 1.95);
    \node[hwblue] at (-1.05,3.0) {orbit};

    % ECI axes
    \draw[->, thick] (O) -- (0,4.1) node[above] {$\vect{I}_z$};
    \draw[->, thick] (O) -- (-2.5,-1.5) node[below left] {$\vect{I}_x$};
    \draw[->, thick] (O) -- (4.6,0.9) node[right] {$\vect{I}_y$};

    % the latitude triangle
    \draw[very thick, hwmagenta] (O) -- (F);
    \node[hwmagenta] at (0.97,-0.73) {$r\cos\phi$};
    \draw[thick, dashed] (F) -- (S) node[pos=0.55, left] {$z_e$};
    \draw[very thick, hwpurple, ->] (O) -- (S) node[midway, above left] {$\vect{r}$};
    \draw (2.086,-0.499) -- (2.086,-0.279) -- (2.3,-0.33);

    % latitude angle
    \draw[->, hwgreen] (-13.45:1.3) arc (-13.45:55.45:1.3);
    \node[hwgreen] at (24:1.7) {$\phi$};

    % spacecraft
    \fill (S) circle (2pt);
    \node[above right] at (S) {spacecraft};
\end{tikzpicture}

\caption{\newtxt{Geocentric latitude $\phi$ at a point on the orbit, from \cref{eq:t2-lonlat}. $z_e = \vect{r}\cdot\vect{I}_z$ is the height above the equatorial plane, and $r\cos\phi = \sqrt{x_e^2+y_e^2}$ is the projection of $\vect{r}$ onto that plane (not to scale).}}
\label{fig:t3-latitude}
\end{figure}

\usertxt{We can expand $\sin\zeta$ with the angle-sum identity for each value of $\omega$ provided to get an expression only in terms of the true anomaly:}

\begin{table}[H]
\centering
\caption{\newtxt{Expansion of $\sin\zeta$ for each value of $\omega$.}}
\label{tab:t3-sinzeta}
\color{newtext}
\begin{tabular}{@{}ccc@{}}
\toprule
$\omega$ & $\sin\zeta = \sin\omega\cos\nu + \cos\omega\sin\nu$ & $\sin\zeta$ \\
\midrule
$90^\circ$ & $(1)\cos\nu + (0)\sin\nu$ & $\cos\nu$ \\
$270^\circ$ & $(-1)\cos\nu + (0)\sin\nu$ & $-\cos\nu$ \\
\bottomrule
\end{tabular}
\end{table}

\noindent Since $i = 63.4^\circ$, $\sin i \approx 0.894$, and $\sin$ is an odd function ($\sin(-x) = -\sin x$):
\begin{equation}
\sin\phi \approx 0.894\,\sin\zeta.
\end{equation}
\usertxt{This means we only need to know the sign of $\sin\zeta$ to know what hemisphere the spacecraft lies in. $\cos(\nu)$ ($\omega = 90$ deg) is negative between 90 deg and 270 deg, and $-\cos(\nu)$ ($\omega = 270$ deg) is negative between 270 deg and 90 deg (passing through 0 deg)}\newtxt{, so those are the ranges where the spacecraft is in the Southern hemisphere.}

The time since periapsis to reach any $\nu$ follows directly from \cref{eq:t1-E0} and Kepler's equation, with no iteration:
\begin{equation}
t(\nu) = \frac{E - e\sin E}{n}, \qquad n = \sqrt{\mu/a^3} = \frac{2\pi}{T}, \qquad E = E(\nu) \text{ from \cref{eq:t1-E0}}.
\label{eq:t3-tnu}
\end{equation}
\end{newpar}

\subsection*{(a)}
\subsection*{(i) Majority hemisphere}

Points with \newtxt{$\phi > 0$ (or $\phi < 0$)} are marked as inside the \newtxt{Northern (or Southern) hemisphere}. Each point stands for one $\Delta t$, so the time in the \newtxt{hemisphere} is (number of points)$\cdot\Delta t$, accurate to about $\pm\Delta t$. \newtxt{The results are listed in \cref{tab:t3a}.}

\subsection*{(ii) Time from the crossing point back to the crossing point}

\begin{newpar}
The points at $\nu$ and $360^\circ - \nu$ have the same latitude (\cref{eq:t3-lat,tab:t3-sinzeta}, since $\cos\nu$ is even). The crossing point occurs at the true anomaly $\nu_c$ where these points also have the same Earth-relative longitude, with $\theta_e$ from \cref{eq:t2-lonlat} at the time $t(\nu)$ of \cref{eq:t3-tnu}:
\begin{equation}
\boxed{\theta_e(\nu_c) = \theta_e(360^\circ-\nu_c)}.
\label{eq:t3-crossing}
\end{equation}
Since $\nu_c$ cannot be found analytically, it is solved numerically as a root-finding problem,
\begin{equation}
\theta_e(\nu_c) - \theta_e(360^\circ-\nu_c) = 0.
\end{equation}
\texttt{scipy.optimize.brentq}~\cite{scipy-brentq} is used because it is a bracketing method that only needs an interval where this difference changes sign, with no derivative. The interval is the pair of neighboring five-minute points in $\vect{N}$ where the sign changes. This gives $\nu_c = 139.31^\circ$, so the two passes are at $\nu = 139.31^\circ$ and $\nu = 220.69^\circ$. These $\nu$'s bound apoapsis, so the time from the crossing point back to it in the majority hemisphere is
\begin{equation}
\boxed{t(220.69^\circ) - t(139.31^\circ) = \SI{8.31}{\hour}}.
\end{equation}
\end{newpar}

\begin{table}[H]
\centering
\caption{\newtxt{Part (a) results.}}
\label{tab:t3a}
\color{newtext}
\begin{tabular}{@{}lrr@{}}
\toprule
Quantity & $\omega=90^\circ$ & $\omega=270^\circ$ \\
\midrule
Majority hemisphere & Southern & Northern \\
Time from crossing point back to crossing point [\si{\hour}] & $8.31$ & $8.31$ \\
\bottomrule
\end{tabular}
\end{table}

\subsection*{(b) $\Omega$ placing apoapsis over Australia and South America ($\omega=90^\circ$)}

\begin{newpar}
As stated before, changing $\Omega$ only shifts the track longitudinally, so the $\Omega$ that places apoapsis ($\nu=180^\circ$, $t = T/2$) over a target longitude $\theta_t$ can be defined as a translation from $\Omega = 0$:
\begin{equation}
\boxed{\Omega = \Big(\theta_t - \theta_e\big|_{\nu=180^\circ,\ \Omega=0}\Big) - 360^\circ\left\lfloor \frac{\theta_t - \theta_e\big|_{\nu=180^\circ,\ \Omega=0}}{360^\circ} \right\rfloor}.
\label{eq:t3-omega-apo}
\end{equation}
For $\omega=90^\circ$, apoapsis is in the Southern hemisphere and $\theta_e = 89.51^\circ$ for $\Omega = 0$. Since $0 \le \Omega \le 180^\circ$ (\cref{eq:t1-given}), any value from \cref{eq:t3-omega-apo} outside that range is replaced by the closest allowed value. The rough reference longitudes in \cref{tab:t3b} give the values of $\Omega$ and the resulting apoapsis longitudes. For South America, this limit leaves apoapsis west of the continent.
\end{newpar}

\begin{table}[H]
\centering
\caption{\newtxt{Part (b) results, $\omega=90^\circ$.}}
\label{tab:t3b}
\color{newtext}
\begin{tabular}{@{}lrrr@{}}
\toprule
Target & Rough reference longitude [\si{\degree}] & $\Omega$ [\si{\degree}] & Apoapsis longitude [\si{\degree}] \\
\midrule
Australia & $135.0$ & $45.4928$ & $135.00$ \\
South America & $-60.0$ & $180.0000$ & $-90.49$ \\
\bottomrule
\end{tabular}
\end{table}

\subsection*{(c) $\Omega$ placing apoapsis over Russia and North America ($\omega=270^\circ$)}

\begin{newpar}
\Cref{eq:t3-omega-apo} applies with $\omega=270^\circ$, where apoapsis is in the Northern hemisphere, and $\theta_e = -90.49^\circ$ for $\Omega = 0$ in this hemisphere. Using the rough reference longitudes in \cref{tab:t3c} gives the values of $\Omega$.
\end{newpar}

\begin{table}[H]
\centering
\caption{\newtxt{Part (c) results, $\omega=270^\circ$.}}
\label{tab:t3c}
\color{newtext}
\begin{tabular}{@{}lrrr@{}}
\toprule
Target & Rough reference longitude [\si{\degree}] & $\Omega$ [\si{\degree}] & Apoapsis longitude [\si{\degree}] \\
\midrule
Russia & $90.0$ & $180.0000$ & $89.51$ \\
North America & $-100.0$ & $0.0000$ & $-90.49$ \\
\bottomrule
\end{tabular}
\end{table}


\clearpage
\subsection*{Full program output}
\programoutput{outputs/main.txt}

\FloatBarrier
